\section{Datasets}
\label{sec:datasets}

The original project evaluated a broad mixture of TU benchmarks, including molecular, vision, and generic graph-classification datasets. For the present submission target, that breadth is not an advantage. We therefore reorganize the experimental evidence around protein-oriented datasets and treat non-protein benchmarks only as supplementary robustness checks. This section describes the topic-facing dataset package used in the main paper.

\subsection{Topic-Facing Dataset Selection}

Our main experiments use three graph classification datasets from the TU Dortmund collection accessed through PyTorch Geometric \cite{fey2019fast}:

\begin{itemize}
    \item \textbf{PROTEINS}: graph-level protein classification
    \item \textbf{DD}: protein structure graph classification on much larger and sparser graphs
    \item \textbf{ENZYMES}: enzyme-oriented multi-class graph classification
\end{itemize}

This selection is deliberate. PROTEINS and DD are the strongest biological anchors already present in our codebase, while ENZYMES improves topic fit with only minimal engineering change because it remains within the same TUDataset interface. By contrast, datasets such as MUTAG, AIDS, and Mutagenicity are retained only as supplementary evidence, and MSRC\_9 is excluded from the topic-facing main text because it is not biologically relevant to the target venue.

\subsection{Dataset Overview}

\paragraph{PROTEINS}
PROTEINS contains graph representations of proteins, where nodes denote secondary structure elements and edges encode neighborhood relations between them. The task is binary graph classification. In the current data export, the dataset contains 1,113 graphs, 2 classes, and 3 input features per node. The average graph has 39.06 nodes and 72.82 edges, although the size variation is substantial.

\paragraph{DD}
DD is a protein structure dataset in which each graph corresponds to a protein and nodes represent amino acids or structure-related units linked by proximity-derived edges. The task is binary classification. DD is the largest and sparsest benchmark in the main package, with 1,178 graphs, 2 classes, and 89 input features. On average, graphs contain 284.32 nodes and 715.66 edges, with individual graphs reaching 5,748 nodes.

\paragraph{ENZYMES}
ENZYMES is a 6-class enzyme classification benchmark with 600 graphs and 3 input features. It is attractive for the present paper because it is protein- and enzyme-oriented while still being easy to integrate into the existing graph-classification pipeline. The average graph has 32.63 nodes and 62.14 edges, and the class distribution is perfectly balanced across the six categories.

\subsection{Dataset Statistics}

Table~\ref{tab:topic_datasets} summarizes the key statistics used in the topic-facing experiments.

\begin{table}[t]
\centering
\caption{Summary of the topic-facing biomolecular graph datasets.}
\label{tab:topic_datasets}
\begin{tabular}{lccccc}
\toprule
Dataset & \# Graphs & \# Classes & Feature Dim. & Avg. Nodes & Avg. Edges \\
\midrule
PROTEINS & 1113 & 2 & 3 & 39.06 & 72.82 \\
DD & 1178 & 2 & 89 & 284.32 & 715.66 \\
ENZYMES & 600 & 6 & 3 & 32.63 & 62.14 \\
\bottomrule
\end{tabular}
\end{table}

\subsection{Data Loading}

All datasets are loaded through the PyTorch Geometric \texttt{TUDataset} interface. Each graph is represented as a \texttt{Data} object containing node features \texttt{x}, graph connectivity \texttt{edge\_index}, graph label \texttt{y}, and batch indices generated during mini-batch loading. If a graph is missing node attributes, the loader replaces them with a constant all-ones feature, ensuring that every benchmark can be processed by the same training code.

No additional graph augmentation, feature normalization, or edge-attribute modeling is applied in the topic-facing protocol. This is a deliberate choice: the goal is to isolate the effect of residual and cross-residual information flow rather than to combine it with a large preprocessing stack.

\subsection{Splitting Protocol}

The present V3 code no longer uses the earlier contiguous fold slicing scheme. Instead, all reported topic-facing results follow a stratified two-stage split:

\begin{itemize}
    \item \textbf{Outer evaluation}: stratified 5-fold cross-validation at the graph level
    \item \textbf{Inner validation}: a stratified random validation subset sampled from the training fold with validation ratio 0.1
\end{itemize}

This change is important. Under weaker settings, several models collapsed toward majority-class behavior, which blurred architectural differences. The stratified outer fold construction and stratified inner validation split now provide a more reliable basis for early stopping and model comparison.

All experiments use a fixed global seed for reproducibility, while the inner train/validation split uses a fold-dependent seed so that every outer fold keeps a deterministic but distinct validation partition.

\subsection{Evaluation Metrics}

The primary evaluation metric is graph classification accuracy on the held-out test fold. For each method and dataset, we report the mean and standard deviation of the best test accuracy over the 5 stratified folds. We also track test loss, best epoch, parameter count, learning-rate history, gradient norms, and embedding statistics, which are used in the analysis scripts and TensorBoard logs.

\subsection{Scope Note}

Although these datasets are more biologically aligned than the generic benchmark mixture used earlier in the project, they still do not constitute a full integrative omics benchmark suite. Accordingly, the claims of this paper are intentionally conservative. We position the method as a biomolecular graph representation framework evaluated on protein-related datasets, not as a complete omics-integration system or a direct solution to novel gene, domain, or motif discovery.
